\begin{problem}{E}{Expected Traversal}{1 second}

\noindent
In algorithm class, we explore a powerful searching algorithm for trees: Depth-First Search (DFS). Consider the following pseudocode for finding a specific target node $U$ in a tree rooted at node $1$:

\begin{verbatim}
cost = 0
found = false

function DFS(node, parent):
    if node == U:
        found = true
        return
        
    children = get_all_neighbors_except(node, parent)
    permute children uniformly at random
    
    for child in children:
        if found == true: 
            return
            
        cost = cost + 1      // Traverse to child
        DFS(child, node)
\end{verbatim}

\noindent
Given a tree with $N$ vertices and a target node $U$, suppose we call \texttt{DFS(1, 0)}. Because the order in which the children are visited is permuted uniformly at random at each step, the final traversal cost (the value of \texttt{cost} when the algorithm terminates) is a random variable. 

\noindent
What is the expected traversal cost to find the target node $U$?

\InputFile

\noindent
The first line contains two integers $N$ and $U$ ($1 \le N \le 10^5$, $1 \le U \le N$) --- the number of vertices in the tree and the target node, respectively.

\noindent
Each of the following $N - 1$ lines contains two integers $u$ and $v$ ($1 \le u, v \le N$), representing an undirected edge between node $u$ and node $v$. It is guaranteed that the given edges form a valid tree.

\OutputFile

\noindent
Print a single integer: the expected traversal cost modulo $10^9 + 7$. 

Formally, let $M = 10^9 + 7$. It can be shown that the exact answer can be expressed as an irreducible fraction $\frac{P}{Q}$, where $P$ and $Q$ are integers and $Q \not\equiv 0 \pmod M$. Output the integer equal to $P \cdot Q^{-1} \bmod M$. In other words, output such an integer $x$ that $0 \le x < M$ and $x \cdot Q \equiv P \pmod M$.

\Examples

\begin{example}
\exmp{1}{
2 2
1 2
}{1
}
\end{example}

\begin{example}
\exmp{2}{
4 2
1 2
1 3
1 4
}{2
}
\end{example}

\begin{example}
\exmp{3}{
3 2
1 2
1 3
}{500000005
}
\end{example}
\Explaination{\upshape\normalsize In the third example, the expected cost is $\frac{3}{2}$. Modulo $10^9 + 7$, is $3 \cdot 500000004 \equiv 500000005$.}

\pagebreak

\end{problem}